\section{Sample variance}\label{sec:sample-var}

Let $X_1,\ldots,X_n$ be i.i.d. copies of a centered variance-one random variable $X$, let $\bar X=n^{-1}\sum_iX_i$, and set
\[
Q_n=\sum_{i=1}^n(X_i-\bar X)^2.
\]
This section applies reflection stability to $Q_n$. We first record the exact consequence of \cref{thm:exact-psd}. For stability, spherical averaging provides an intermediate transform quantity, which will be bounded in terms of the usual Kolmogorov distance in Section~\ref{subsec:sample-kolmogorov}.

\subsection{The exact consequence}

The exact sample-variance consequence of the quadratic-form theorem is the following.
\begin{corollary}[Exact sample-variance characterization under dominance]\label[corollary]{cor:exact-sample-var}
Let $n\ge2$ and let $X_1,\ldots,X_n$ be i.i.d., centered, and variance one. If
\[
M_X(s)M_X(-s)\ge e^{s^2}
\]
for all $s$ in some neighborhood of the origin, and
\[
\sum_{i=1}^n(X_i-\bar X)^2\sim\chi^2_{n-1},
\]
then $X_i\sim\Normal(0,1)$.
\end{corollary}

\begin{proof}
Take $A=I_n-n^{-1}\mathbf1\mathbf1^{\mathsf T}$ in \cref{thm:exact-psd}. This is a rank-$(n-1)$ orthogonal projection with every diagonal entry equal to $(n-1)/n>0$, and $X^{\mathsf T}AX=Q_n$.
\end{proof}

As in \cref{thm:exact-psd}, the chi-square law implies MGF existence for every coordinate. When $n=2$, dominance is not needed for the exact conclusion: $Q_2=(X_1-X_2)^2/2$, and the symmetry of $X_1-X_2$ identifies its law as $\Normal(0,2)$. Cram\'er's theorem then gives Gaussianity of $X$.

\subsection{From spherical discrepancy to reflected defect}

For $r\ge1$, let $V$ be uniform on $S^{r-1}$. Recall the spherical transform
\[
\Psi_r(z):=\E_V e^{zV_1},
\]
and define, for a nonnegative random variable $Q$ with finite displayed expectations,
\[
\mathfrak D_{r,\tau}(Q):=
\sup_{|t|\le\tau}
\left|e^{-t^2/2}\E\Psi_r(t\sqrt Q)-1\right|.
\]
This quantity vanishes for $Q\sim\chi_r^2$. It is a transform discrepancy, not a probability metric. Its role is to connect the quadratic form to an average of the one-dimensional reflected defect. The next subsection bounds it by Kolmogorov distance under an exponential envelope.

The quantitative statement based on the spherical observable is the next theorem.
\begin{theorem}[Sample-variance stability under spherical Laplace discrepancy]\label{thm:sample-var}
Fix $n\ge3$, $\tau,K>0$, and $0<\rho'<\tau\sqrt{(n-1)/n}$. There are constants $C<\infty$ and $\varepsilon_0\in(0,e^{-e})$, depending only on $(n,\tau,K,\rho')$, such that the following holds. Let $X_1,\ldots,X_n$ be i.i.d. copies of a centered variance-one $X$ with $\E e^{2\tau|X|}\le K$, and assume $R_X(s)\ge0$ for $|s|\le\tau$. If
\[
Q_n=\sum_{i=1}^n(X_i-\bar X)^2,
\qquad
\mathfrak D_{n-1,\tau}(Q_n)\le\varepsilon,\qquad 0<\varepsilon\le\varepsilon_0,
\]
then
\[
\dK(X,\Normal(0,1))
\le C\sqrt{\frac{\log\log(1/\varepsilon)}{\log(1/\varepsilon)}}.
\]
\end{theorem}

\begin{proof}
Let $r=n-1$, let $P=I_n-n^{-1}\mathbf1\mathbf1^{\mathsf T}$, and choose $B\in\R^{r\times n}$ with $B^{\mathsf T}B=P$ and $BB^{\mathsf T}=I_r$. If $V$ is uniform on $S^{r-1}$, then every column $b_i$ of $B$ has norm
\[
q:=\|b_i\|=\sqrt{\frac{n-1}{n}},
\]
and $\langle b_i,V\rangle$ has the same distribution as $qV_1$. Put $a=\tau q$ and define
\[
D_X(s):=\frac12R_X(s)
 =\frac12\log\frac{M_X(s)M_X(-s)}{e^{s^2}}.
\]
Under dominance, $D_X\ge0$ on $[-\tau,\tau]$.

The exact spherical transform gives, for $|t|\le\tau$,
\begin{equation}\label{eq:spherical-discrepancy-identity}
e^{-t^2/2}\E\Psi_r(t\sqrt{Q_n})
=\E_V\exp\left(\sum_{i=1}^n\left[\log M_X(t\langle b_i,V\rangle)-\frac{t^2\langle b_i,V\rangle^2}{2}\right]\right).
\end{equation}
Let $H_t(V)$ denote the exponent on the right. The identity uses $BB^{\mathsf T}=I_r$, which gives $\sum_i\langle b_i,V\rangle^2=1$. Each $\langle b_i,V\rangle$ has the law of $qV_1$, and symmetry gives
\[
\E_V H_\tau(V)=n\E D_X(aV_1)\ge0.
\]
Jensen's inequality and the definition of the discrepancy now give
\[
1\le\exp\{n\E D_X(aV_1)\}
\le e^{-\tau^2/2}\E\Psi_r(\tau\sqrt{Q_n})
\le1+\mathfrak D_{n-1,\tau}(Q_n).
\]
Taking logarithms yields
\begin{equation}\label{eq:integrated-D}
n\E D_X(aV_1)\le\log\bigl(1+\mathfrak D_{n-1,\tau}(Q_n)\bigr).
\end{equation}
We next turn this averaged estimate into a bound on an interior interval.

For the pointwise transfer, the density of $V_1$ is
\[
f_n(v)=
\frac{\Gamma((n-1)/2)}{\sqrt\pi\,\Gamma((n-2)/2)}
(1-v^2)^{(n-4)/2},
\qquad -1<v<1.
\]
Fix $0<\rho'<a$ and choose $0<h<a-\rho'$. On $[-(\rho'+h),\rho'+h]$ the scaled density of $aV_1$ is bounded below by
\[
m_h:=\frac1a\frac{\Gamma((n-1)/2)}{\sqrt\pi\,\Gamma((n-2)/2)}
\left(1-\left(\frac{\rho'+h}{a}\right)^2\right)^{\max(n-4,0)/2}.
\]
The exponential envelope implies that $D_X$ is Lipschitz there. One may take
\[
L_h:=\frac{K}{e\tau}+\rho'+h,
\]
because $M_X(s)\ge1$ and $|x|e^{\tau|x|}\le(e\tau)^{-1}e^{2\tau|x|}$. If $M_*:=\sup_{|s|\le\rho'}D_X(s)$, choose $s_0\in[-\rho',\rho']$ with $D_X(s_0)=M_*$. The function $D_X$ is even, nonnegative, and satisfies $D_X(0)=0$. Since it is $L_h$-Lipschitz on the enlarged interval,
\[
D_X(s)\ge\bigl(M_*-L_h|s-s_0|\bigr)_+.
\]
If $M_*=0$, the desired pointwise bound is immediate. Otherwise $|s_0|\ge M_*/L_h$ because $D_X(0)=0$. Put $\ell=\min\{h,M_*/L_h\}$. The intervals
\[
[s_0-\ell,s_0+\ell],\qquad [-s_0-\ell,-s_0+\ell]
\]
lie in the enlarged interval and have disjoint interiors, since $\ell\le|s_0|$. On each, integrate the corresponding tent. If $M_*\le L_hh$, its positive part has half-width $M_*/L_h$; if $M_*>L_hh$, restrict the tent to half-width $h$. Integrating against the density lower bound gives
\[
2m_h I_{L_h,h}(M_*)
 \le \frac1n\log\bigl(1+\mathfrak D_{n-1,\tau}(Q_n)\bigr),
\]
where
\[
I_{L,h}(M):=
\begin{cases}
M^2/L,& M\le Lh,\\
2hM-Lh^2,& M>Lh.
\end{cases}
\]
In both cases $I_{L,h}(M)=\int_{-h}^h(M-L|u|)_+\,du$. Writing $G(y;L,h)=\sqrt{Ly}$ when $y\le Lh^2$ and $G(y;L,h)=y/(2h)+Lh/2$ otherwise, we obtain
\begin{equation}\label{eq:pointwise-D}
\sup_{|s|\le\rho'}D_X(s)
\le\inf_{0<h<a-\rho'}G\left(\frac{\log(1+\mathfrak D_{n-1,\tau}(Q_n))}{2nm_h};L_h,h\right).
\end{equation}
In particular, for fixed $n,\tau,K,\rho'$, the right-hand side is $O(\sqrt{\mathfrak D_{n-1,\tau}(Q_n)})$ as the discrepancy tends to zero. The strict interior restriction is used by this density-lower-bound argument when $n\ge5$, because the spherical density vanishes at the endpoints; an endpoint result would require a different local estimate.

Applying \eqref{eq:pointwise-D} to the sample-variance hypothesis gives
\[
\Delta_{\rho'}(X)
\le C_1\sqrt{\varepsilon},
\]
for all sufficiently small $\varepsilon$, where $C_1=C_1(n,\tau,K,\rho')$; here we used $R_X=2D_X$. Set $\Delta=C_1\sqrt{\varepsilon}$. Then
\[
\log(1/\Delta)=\tfrac12\log(1/\varepsilon)-\log C_1
\asymp \log(1/\varepsilon),
\qquad
\log\log(1/\Delta)=\log\log(1/\varepsilon)+O(1).
\]
Theorem~\ref{thm:sharp-upper}, with radius $\rho'$, yields the asserted rate. The exponential envelope was used to obtain the Lipschitz bound that transfers the spherical average to a pointwise defect.
\end{proof}

\subsection{Transfer from Kolmogorov distance}\label{subsec:sample-kolmogorov}

We now control the intermediate transform discrepancy by Kolmogorov distance. Let $H_r\sim\chi_r^2$.

\begin{lemma}[Kolmogorov control of the spherical observable]\label[lemma]{lem:kolmogorov-to-spherical}
Fix $r\ge1$ and $\tau>0$. Let $Q\ge0$ and suppose
\[
\E e^{2\tau\sqrt Q}\le A.
\]
Then, for
\[
\delta:=\dK(Q,H_r)\le1,
\]
one has
\[
\mathfrak D_{r,\tau}(Q)\le C_{r,\tau,A}\sqrt\delta.
\]
\end{lemma}

\begin{proof}
If $\delta=0$, then $Q$ and $H_r$ have the same law and the discrepancy is zero. Suppose $0<\delta\le1$. For $|t|\le\tau$, define
\[
g_t(q):=e^{-t^2/2}\Psi_r(t\sqrt q),\qquad q\ge0.
\]
Because $V_1$ is symmetric,
\[
\Psi_r(t\sqrt q)=\E_V\cosh(t\sqrt q\,V_1),
\]
so $g_t$ is nonnegative and nondecreasing in $q$. Also
\[
0\le g_t(q)\le e^{|t|\sqrt q}\le e^{\tau\sqrt q}.
\]
If $H_r=\|G\|^2$ with $G\sim\Normal(0,I_r)$, spherical symmetry gives
\[
\E g_t(H_r)=e^{-t^2/2}\E e^{tG_1}=1.
\]
Moreover,
\[
\sup_{|t|\le\tau}\E g_t(Q)^2\le A,
\qquad
\sup_{|t|\le\tau}\E g_t(H_r)^2
\le \E e^{2\tau\sqrt{H_r}}=:A_{r,\tau}<\infty.
\]
For $T\ge1$, the function $g_t\wedge T$ is bounded and nondecreasing, with total variation at most $T$. Integration by parts against the difference of the distribution functions therefore gives
\[
\left|\E(g_t(Q)\wedge T)-\E(g_t(H_r)\wedge T)\right|
\le T\delta.
\]
The tails satisfy
\[
\E(g_t(Q)-T)_+\le\frac{\E g_t(Q)^2}{T}\le\frac{A}{T},
\]
and likewise for $H_r$. Hence, uniformly in $|t|\le\tau$,
\[
|\E g_t(Q)-1|
\le T\delta+\frac{A+A_{r,\tau}}{T}.
\]
Taking $T=\delta^{-1/2}$ proves the claim.
\end{proof}

For the sample variance, the exponential envelope already assumed in \cref{thm:sample-var} provides the hypothesis of \cref{lem:kolmogorov-to-spherical}. Indeed, since $Q_n=X^{\mathsf T}PX\le\sum_iX_i^2$,
\[
\sqrt{Q_n}\le\left(\sum_iX_i^2\right)^{1/2}\le\sum_i|X_i|,
\]
and independence gives
\begin{equation}\label{eq:sample-var-radial-envelope}
\E e^{2\tau\sqrt{Q_n}}
\le\prod_{i=1}^n\E e^{2\tau|X_i|}
\le K^n.
\end{equation}
We therefore obtain the following standard-distance consequence.

\begin{corollary}[Sample-variance stability in Kolmogorov distance]\label[corollary]{cor:sample-var-kolmogorov}
Fix $n\ge3$, $\tau,K>0$, and $0<\rho'<\tau\sqrt{(n-1)/n}$. There exist constants $C<\infty$ and $\delta_0\in(0,e^{-e})$, depending only on $(n,\tau,K,\rho')$, such that the following holds. Let $X_1,\ldots,X_n$ be i.i.d. copies of a centered variance-one $X$ satisfying
\[
\E e^{2\tau|X|}\le K,
\qquad
R_X(s)\ge0\quad(|s|\le\tau).
\]
If
\[
0<\delta:=\dK(Q_n,\chi^2_{n-1})\le\delta_0,
\]
then
\[
\dK(X,\Normal(0,1))
\le C\sqrt{\frac{\log\log(1/\delta)}{\log(1/\delta)}}.
\]
If $\delta=0$, then $X\sim\Normal(0,1)$.
\end{corollary}

\begin{proof}
The case $\delta=0$ follows from \cref{cor:exact-sample-var}. For $\delta>0$, by \eqref{eq:sample-var-radial-envelope} and \cref{lem:kolmogorov-to-spherical},
\[
\mathfrak D_{n-1,\tau}(Q_n)\le C_0\sqrt\delta.
\]
Apply \cref{thm:sample-var} with $\varepsilon=C_0\sqrt\delta$. Since
\[
\log(1/\varepsilon)=\tfrac12\log(1/\delta)+O(1),
\qquad
\log\log(1/\varepsilon)=\log\log(1/\delta)+O(1),
\]
the asserted rate follows after changing the constant.
\end{proof}

\subsection{Scope and further questions}

The application uses dominance to recover pointwise information from a nonnegative spherical average. The reflection theorem itself has no such sign assumption. The restriction $n\ge3$ in the quantitative argument allows us to use the density of a spherical coordinate; the exact $n=2$ case was treated separately above.

The sharpness result in Section~\ref{sec:sharpness} concerns the reflected-MGF problem and does not establish sharpness of \cref{cor:sample-var-kolmogorov}. Determining the optimal sample-variance modulus remains a separate question. It would also be useful to weaken the exponential envelope used to control the unbounded spherical transform, or to develop a quantitative counterpart of \cref{thm:exact-psd} for non-identically distributed coordinates.
